\subsection{把穿孔权重用于能量计算}
分子从小孔逸出是一个条件抽样过程：容器内随机选取一个分子，与随机选取一个实际穿孔事件，不具有同样的速度分布。前面的通量积分中 $v_z$ 权重既影响穿孔数，也影响每个逸出分子平均带走的能量。

\begin{origexamplebox}{小孔逸出为何使绝热容器中的气体冷却}\label{ch18:ex:teaching-effusion}
刚性绝热容器中有稀薄单原子理想气体，粒子数和温度初值为 $N_0,T_0$。通过足够小的孔缓慢向真空逸出；逸出之间气体有足够时间恢复麦克斯韦分布。忽略孔边作用和回流。求一个逸出分子的平均动能，以及容器内剩余气体的温度随 $N$ 的关系。
\end{origexamplebox}
\begin{xesolution}
取孔面外法线为 $z$ 轴。麦克斯韦分布三个速度分量独立；穿孔权重只含 $v_z>0$，故两个切向分量仍各贡献 $k_BT/2$。令 $a=m/(2k_BT)$，法向分量的加权二阶矩为
\[
 \langle v_z^2\rangle_{\rm hit}
 =\frac{\int_0^\infty v_z^3e^{-av_z^2}\dd v_z}
        {\int_0^\infty v_ze^{-av_z^2}\dd v_z}
 =\frac{1/(2a^2)}{1/(2a)}=\frac1a.
\]
两积分用 $u=av_z^2$ 即可得到。法向平均动能为 $m/(2a)=k_BT$，因此每个逸出分子平均携带
\[
 \overline E_{\rm out}=\frac12k_BT+\frac12k_BT+k_BT=2k_BT,
\]
大于容器内分子的 $3k_BT/2$。快分子更易逸出，剩余分子因而冷却。
容器边界固定且不传热，但有物质携能离开，不能写成 $\dd U=0$。在 $\dd N<0$ 的一步中，流出的粒子数为 $-\dd N$，能量变化为 $\dd U=2k_BT\dd N$。而
\[
 \dd U=\frac32k_B(N\dd T+T\dd N).
\]
两式相等得到 $3N\dd T=T\dd N$，积分为
\[
 \ln\frac{T}{T_0}=\frac13\ln\frac{N}{N_0},\qquad
 T=T_0\left(\frac{N}{N_0}\right)^{1/3}.
\]
$N$ 减小时 $T$ 减小，初值和符号均相符。这不是封闭气体的可逆绝热膨胀，不能套 $TV^{\gamma-1}=\mathrm{const}$；本题体积没有变化，变的是粒子数。若漏气过快、孔尺度超过平均自由程或来不及再平衡，穿孔抽样及此冷却律都须重新判断。
\end{xesolution}
